\chapter{Development of the Interactive Selection Tool}
\label{ch:selection_tool}

To translate the 24-case CFD Design of Experiments into an accessible engineering utility, an interactive software application is planned using the Python \texttt{Streamlit} framework \cite{streamlit_docs}. This chapter describes the software architecture, data aggregation logic, and decision-making algorithm intended to help users compare candidate valve configurations within the verified simulation range.

\section{Software Architecture and Data Aggregation}
\label{sec:software_architecture}
The proposed selection tool separates offline data reduction and quality checks from user interaction. The supplied archive has twelve hydrogen post-processing datasets but no nitrogen results; therefore, it cannot yet support an independently verified choice between the two gases:
\begin{enumerate}
    \item \textbf{Converged State Extraction:} Using the parsing functions established in \texttt{generate\_reports.py}, the accepted steady-state values (after residual, force, pressure and continuity checks) for pressure drop $\Delta P$, axial obturator force $F_x$, and volume flux $Q$ are extracted from all 24 \texttt{postProcessing/} directories and compiled into a structured lookup table. %
    \item \textbf{Dimensional Scaling Engine:} Kinematic pressure differences would be converted into pascals using the selected reference density. The new hydrogen force files have already been dimensionalized by the active \texttt{rhoInf = 0.0838} value, while nitrogen case dictionaries specify \texttt{rhoInf = 1.165}. Therefore, no further density multiplication should be applied to the output forces. Extrapolation to different tank pressures and temperatures would require a separate assessment of changes in density, viscosity and possible compressibility \cite{white2016,nist_refprop}.
\end{enumerate}

\section{User Interface and Valve Recommendation Logic}
\label{sec:recommendation_logic}
Through the web-based graphical user interface (GUI), the user specifies four operational constraints:
\begin{itemize}
    \item Working gas selection ($\text{H}_2$ or $\text{N}_2$). %
    \item Operating tank pressure and temperature.
    \item Nominal required flow rate ($Q_{\text{nom}}$ or $\dot{m}_{\text{nom}}$) and maximum allowable pressure drop ($\Delta p_{\text{max}}$) during normal operation.
    \item Target excess flow trip threshold ($Q_{\text{trip}}$) and pre-load spring force ($F_{\text{spring}}$) required to actuate the obturator during a line rupture.
\end{itemize}

Once a verified dataset is available, the tool may interpolate CFD responses over the simulated inlet-velocity range separately for each gas and obturator geometry (\texttt{3Fori\_1p5mm}, \texttt{3Fori\_3mm}, \texttt{6Fori\_1p5mm}, and \texttt{6Fori\_3mm}). %
The proposed logic would exclude geometries that exceed $\Delta p_{\text{max}}$ at nominal flow or fail to generate $F_x \ge F_{\text{spring}}$ at a specified comparison flow rate, subject to consistency of the force direction and a separately validated closure criterion. A preliminary implementation would display comparative Plotly performance curves, identify its interpolation limits and avoid claiming experimentally validated trip thresholds \cite{plotly_docs, streamlit_docs}.
 
\section{Available Dataset and Implementation Limits}
\label{sec:selection_coverage}

The current archive provides preliminary hydrogen response data for the complete $3\times4$ inlet-velocity--obturator matrix, but no nitrogen post-processing results. A user interface should distinguish available, verified and missing data explicitly. Predictions should not be generated for missing nitrogen conditions, and extrapolation beyond the simulated velocities should be disabled or clearly labeled as unsupported.

The source case files contain reference-state viscosity and density values approximately corresponding to $20\,^{\circ}\mathrm{C}$ and $1\,\mathrm{atm}$. They do not demonstrate valve performance at the separately proposed $437\,\mathrm{bar}$, $-40\,^{\circ}\mathrm{C}$ operating condition. Applying the comparison tool to that state would require suitable real-fluid properties, consideration of compressibility and, where needed, additional CFD runs. Finally, steady fixed-obturator hydrodynamic forces cannot on their own validate the transient closing or tripping behavior of the valve.
